\BeleskeID{03-s013}
% element: 03-s013:e01
% element: 03-s013:e02
% element: 03-s013:e03
% element: 03-s013:d01

Multiplekser vodi više podatkovnih ulaza ka jednom izlazu, pa se po odnosu njihovog broja može posmatrati kao posebna koderska mreža. On ipak ne koduje indeks aktivnog ulaza: adresu zadaju selektori, a izlaz prenosi \emph{vrednost} izabranog podatka. U odnosu na demultiplekser obrnut je smer izbora, više-prema-jedan umesto jedan-prema-više.

Svaka podatkovna grana uslovljena je svojim mintermom selektora. Za MUX 4/1 važi
\[
 Z=I_0\overline{S_1}\overline{S_0}
  +I_1\overline{S_1}S_0
  +I_2S_1\overline{S_0}+I_3S_1S_0.
\]
Za svaku ustaljenu adresu samo jedan minterm je jedinica, pa ILI kolo ne meša više podataka: propušta baš \(I_{2S_1+S_0}\). Na primer, adresa 01 bira \(I_1\), bez obzira na vrednosti ostala tri podatkovna ulaza.
